% SPDX-License-Identifier: Apache-2.0
% covers: Raster
\datasheet{Raster}{Razboj's rasteriser: it reads a display list over AXI and writes one pixel a cycle.}

\dsfacts{\code{gpu/razboj/src/raster.rs}}{\code{razboj::raster::Raster}}%
{\code{//docs:razboj}, \code{//docs:noc}}

\subsection*{Function}

\code{Raster} finds its work in memory and draws it there. It polls a
count at \code{CTRL} until the count is not zero. It then reads the
six words of each instruction of the display list at \code{DL}, one
read at a time, and walks the instruction's box one pixel per cycle.
It writes the instruction's colour at every pixel inside the
primitive: every pixel of a clear or a rectangle, and every pixel of a
triangle at which no edge function is negative.

It is an AXI host written as hardware. It issues bursts on
\code{issue}, sends a pixel's beat on \code{wbeat}, and hands
identifiers back on \code{release}. It receives \code{grant},
\code{done} and \code{rdata} from its tracker. \code{idle} goes high
when the list is drawn and no write is in flight.

\subsection*{Parameters}

\begin{center}\footnotesize
\begin{tabular}{@{}lp{0.7\textwidth}@{}}
\toprule
Parameter & Meaning \\
\midrule
\code{A} & The width of the link's address, in bits. \\
\code{I} & The width of the AXI identifier, in bits. \\
\code{LOGW} & The screen is \code{1 << LOGW} pixels wide. \\
\code{H} & The screen is \code{H} pixels high. \\
\code{BASE} & The byte address of the framebuffer's first word. \\
\code{DL} & The byte address of the display list's first word. \\
\code{CTRL} & The byte address of the count. \\
\bottomrule
\end{tabular}
\end{center}

The tables use \code{Raster<16, 2, 4, 16, 0, 0x400, 0x600>}, the
rasteriser of \code{//gpu/razboj:trace}: a sixteen by sixteen screen at
\code{0x0}, the list at \code{0x400} and the count at \code{0x600}. In
\code{soc} it is \code{Raster<32, 2, 7, 96, 0x8000, 0x2000, 0x4000>}.

\subsection*{Address map}

\begin{center}\footnotesize
\begin{tabular}{@{}lp{0.6\textwidth}@{}}
\toprule
Address & What the rasteriser does there \\
\midrule
\code{CTRL} & Reads the count while polling. \\
\code{DL + (i << 5) + (k << 2)} & Reads word \code{k}, 0 to 5, of
instruction \code{i}. \\
\code{BASE + (((y << LOGW) + x) << 2)} & Writes pixel \code{x},
\code{y}. \\
\bottomrule
\end{tabular}
\end{center}

An instruction's words are in the format \code{gpu/razboj/src/dl.rs} states.
Words 6 and 7 of an instruction are never read.

\begin{center}\footnotesize
\begin{tabular}{@{}lp{0.42\textwidth}p{0.3\textwidth}@{}}
\toprule
Word & Low field & High field \\
\midrule
0 & bits 1 to 0: kind, 0 clear, 1 rectangle, 2 triangle & bits 25 to
2: colour \\
1 & bits 9 to 0: \code{x0} & bits 25 to 16: \code{y0} \\
2 & bits 9 to 0: \code{x1} & bits 25 to 16: \code{y1} \\
3 & bits 11 to 0: \code{ax} & bits 27 to 16: \code{ay} \\
4 & bits 11 to 0: \code{bx} & bits 27 to 16: \code{by} \\
5 & bits 11 to 0: \code{cx} & bits 27 to 16: \code{cy} \\
\bottomrule
\end{tabular}
\end{center}

\dstables{Raster}

\subsection*{Behaviour}

\begin{itemize}
\item \code{phase} is 0 while polling, 1 while fetching, 2 while
drawing and 3 when done.
\item While polling or fetching, the rasteriser issues one read when
no read is out and \code{issue} has room. \code{waiting} is set until
the read's beat arrives.
\item A count that is not zero starts the fetch at instruction 0, word
0, and its low eight bits go into \code{left}. A count of zero means
poll again.
\item During the fetch each word's fields are latched as the word
lands. A kind field of 3 is taken as a triangle.
\item When word 5 lands, the walk is set up in that cycle. The third
vertex comes straight from the bus. The six products of the edge
functions are computed then, and nowhere else.
\item A clear's box is the whole screen, from its own type. A
rectangle's and a triangle's box is \code{x0}, \code{y0} to
\code{x1}, \code{y1}. Vertices are sign-extended from twelve bits.
\item Each edge is kept as its value at this pixel, its value at the
start of this row and its two steps, in 32-bit two's complement. A
step to the right adds the column step; a new row reloads from the row
value plus the row step.
\item \code{hit} is high when the kind is not a triangle, or when all
three edge values have bit 31 clear.
\item A pixel with \code{hit} high is written when both \code{issue}
and \code{wbeat} have room. The walk steps when the pixel needs no
write or its write went out, so backpressure holds the walk.
\item A pixel's write is a burst of one beat, size 2, \code{INCR},
with the colour zero-extended to 32 bits, strobe \code{0xf} and
\code{last} set.
\item At the end of the box the rasteriser fetches the next instruction
while \code{left} is above 1, and otherwise goes to phase 3.
\item \code{inflight} counts writes issued whose response has not come
back.
\item When a write response and a read beat arrive together, the
write's identifier goes back first and the read's waits a cycle.
Grants are taken and dropped.
\end{itemize}

\subsection*{Verification}

\begin{itemize}
\item \code{//gpu/razboj:razboj\_test} renders scenes on the rasteriser through
an \code{AxiHost}, an \code{AxiPer} and \code{Fb}, and compares every
pixel with the model of \code{gpu/razboj/src/model.rs}:
\code{a\_scene\_of\_every\_kind\_agrees\_with\_the\_model},
\code{a\_clear\_says\_only\_its\_colour},
\code{a\_triangle\_is\_the\_same\_whichever\_way\_it\_is\_wound},
\code{a\_triangle\_hanging\_off\_the\_screen\_is\_clipped},
\code{a\_single\_pixel\_and\_an\_empty\_box}, and
\code{pseudorandom\_scenes\_agree\_with\_the\_model}, twelve scenes of
up to four random rectangles and triangles.
\item \code{an\_instruction\_survives\_the\_format} in the same target
checks the display list's encoder against its decoder.
\item The build simulates the netlist of \code{Raster<16, 2, 4, 16, 0,
0x400, 0x600>} against the trace of \code{//gpu/razboj:trace} under nvc and
Verilator: \code{//docs:razboj\_sim\_raster\_raster\_tb\_test} and
\code{//docs:razboj\_vsim\_raster\_test}.
\item \code{//soc:demo\_test} runs
\code{a\_core\_draws\_a\_solid\_and\_a\_gpu\_fills\_it}. Vreteno writes
a display list on the lattice, the rasteriser draws it, and every pixel
must match the model's picture of that list.
\end{itemize}

\subsection*{Used by}

\code{razboj::sim::run} and \code{razboj::sim::run\_list}, and so Razboj's
tests, \code{//gpu/razboj:trace} and \code{//gpu/razboj:demo}. The system in
\code{soc}, where the rasteriser is the host at corner 1,0 of the
lattice.

\subsection*{Limits}

\code{//docs:razboj} lists them. It writes one pixel per cycle at most. It
has no depth buffer, so primitives are drawn in list order. It has no
interpolation, so a triangle is one colour, and no texture. It clips
no geometry, only boxes, and the assembler \code{Op::encode} does that
on the host. It has one read out at a time.

The code adds these. It has no reset port, and phase 3 is final, so it
draws one list. \code{left} is eight bits. The pixel's offset from
\code{BASE} is formed at sixteen bits and then widened to \code{A}
bits. A colour is 24 bits.
